%% Extended Abstract for AI4SciSci 2026 (CEUR-WS style, anonymized for double-blind review)
\documentclass[
% twocolumn,
% hf,
]{ceurart}

\sloppy

\usepackage{listings}
\lstset{breaklines=true}

\begin{document}

%% Rights management information.
\copyrightyear{2026}
\copyrightclause{Copyright for this paper by its authors.
  Use permitted under Creative Commons License Attribution 4.0
  International (CC BY 4.0).}

%% Conference information
\conference{AI4SciSci 2026: The 3rd International Workshop on Artificial Intelligence
  for the Science of Science, October 13, 2026, Dallas, TX, USA}

%% Title
\title{Accelerating the Science-to-Technology Transition: How Generative AI Shapes
  the Speed and Breakthroughness of Corporate Innovation}

%% ANONYMIZED FOR DOUBLE-BLIND REVIEW
\author[1]{Anonymous Author(s)}
\address[1]{Anonymous Institution}

\begin{abstract}
Corporate R\&D draws on academic literature to develop commercial technologies, yet the transition from published science to patented invention has historically taken years. Did the post-2023 adoption of generative AI (GenAI) shorten the lag between academic publication and corporate patent citation, and does GenAI-era R\&D reorient what science patents build on? This extended abstract reports preliminary results from an ongoing study of these questions. Using 429,111 patent--paper citation links from the Lens covering US filings from 2018 to 2025, we estimate a difference-in-differences design: technology classes with high GenAI exposure (CPC G06N, G16B, G16C, C40B) form the treated group, mature mechanical classes (F16B, F16H, B65D) form the control group, and the treatment period begins in 2023. We report three preliminary findings. First, the citation-weighted lag between a cited paper's publication and the citing patent's filing \emph{lengthened} in GenAI-exposed classes relative to controls after 2023 (about +0.20 log points, robust across progressively stricter truncation guards), while the median lag of the typical patent was unchanged: the effect is a composition effect, not a slowdown of the average patent. Second, the mechanism is a reorientation of citation portfolios toward the older canonical stock: the share of citations to papers three years old or younger fell from 19.7\% to 11.9\% in treated classes (a patent-level difference-in-differences of $-4.9$ percentage points), with mid-2010s deep-learning classics absorbing the difference. Third, a placebo estimation with a false 2017 treatment on 2015--2019 data shows no effect, and fresh-science intake in treated classes was \emph{rising} before 2023; the post-2023 decline is a trend reversal rather than field aging. Contrary to the acceleration narrative, the first patent-record evidence suggests GenAI-era invention is consolidating around foundational science rather than speeding up the uptake of new science.
\end{abstract}

\begin{keywords}
  generative AI \sep
  science of science \sep
  science-to-technology lag \sep
  patents \sep
  breakthrough innovation \sep
  difference-in-differences \sep
  semantic homogenization
\end{keywords}

\maketitle

\section{Introduction and Motivation}

Corporate research and development depends on academic science. Firms read the literature, absorb new findings, and translate them into commercial technologies that surface as patents. This translation is slow. Prior work has documented a science-to-technology lag that typically spans years between the publication of a scientific result and its citation in a patent \cite{ahmadpoor2017dual, marx2020reliance}. In biomedicine, patent citations trail scientific citations by a median of six years for the majority of papers \cite{ke2018comparing}, and the time to first patent citation averages roughly a decade, varying systematically with paper characteristics such as basicness and novelty \cite{ke2020basicness}.

The environment changed after 2023, when generative AI (GenAI) gave R\&D teams tools to synthesize scientific knowledge, automate code, and screen literature at scale. Whether this improves corporate innovation is open: GenAI may help firms connect distant scientific fields, or rival firms relying on the same foundation models may converge on faster but less distinctive, less disruptive patents \cite{funk2017dynamic, park2023papers}.

\section{Research Questions}

\begin{itemize}
\item \textbf{RQ1 (Speed).} To what extent has GenAI adoption after 2023 shortened the lag between academic publication and corporate patent citation?
\item \textbf{RQ2 (Homogenization, descriptive).} Do patent abstracts in GenAI-exposed classes show signs of semantic convergence across competing firms after 2023?
\end{itemize}

\section{Data and Methods}

\textbf{Data.} We draw patent records and their citations to scholarly works from the Lens,\footnote{\url{https://www.lens.org}} which links patents to the papers they cite with dated records on both sides, complemented by OpenAlex for paper-level attributes \cite{priem2022openalex}. The sample covers filings from 2018 through 2025 in seven CPC classes selected from the official CPC definitions maintained by the USPTO and EPO: a GenAI-exposed treated group (G06N, artificial intelligence and machine learning; G16B, bioinformatics; G16C, computational chemistry and materials; C40B, combinatorial chemistry) and a low-exposure mechanical control group (F16B, fastening devices; F16H, gearing; B65D, containers and packaging).

\textbf{Identification.} We estimate difference-in-differences models on the citation links: the log lag (days between paper publication and patent filing) and the patent-level share of citations to recent papers ($\le$3 years old) are regressed on the interaction of treated class and post-2023 filing, with class and filing-year fixed effects and standard errors clustered by class--year. Field-wide shocks common to both groups, such as examination speed-ups or drifting citation practices, are absorbed by the design.

\textbf{Results (preliminary).} The treated$\times$post coefficient on the pair-level log lag is $+0.20$ ($p<0.001$) under a uniform ten-year citation window, and strengthens rather than attenuates as truncation guards tighten; at the patent level the median-lag effect is null, locating the change in citation composition rather than in the typical patent's speed. The patent-level share of fresh science falls by $4.9$ percentage points in treated classes ($p=0.001$; from 19.7\% to 11.9\%), and the modal publication years cited by post-2023 treated patents are 2015--2018. Because paper characteristics such as basicness and novelty shape citation speed \cite{ke2020basicness}, paper-level controls are being added next; the composition finding is robust to restricting to filings through mid-2024.

\textbf{Robustness.} A placebo estimation on 2015--2019 with a false 2017 treatment date yields no effect on either margin (lag: $p=0.46$; and fresh-science share was rising in treated classes pre-2023, falsifying a field-maturation account). For RQ2, term co-occurrence maps of treated-class abstracts before and after 2023 in VOSviewer \cite{vaneck2010vosviewer} provide descriptive evidence. We acknowledge the documented limitations of direct citation linkage \cite{denter2025novelty}; we retain citation links because our question is temporal, and direct citations are the only linkage yielding two dated events between which a lag can be measured.

\section{Contributions}

The study speaks to two audiences. For corporate strategy, the evidence cautions R\&D leaders that GenAI-era invention is, so far, consolidating around canonical science rather than accelerating fresh-science uptake; absorptive-capacity investments in frontier literature remain a differentiator. For the science of science, the results connect the GenAI shock to the debate on declining disruptiveness \cite{park2023papers}: a measurable reallocation of patent citations from the frontier to the canon, identified against mechanical control classes with a clean placebo. All queries, code, cleaning rules, results, and the cleaned dataset are available in the anonymized replication archive.\footnote{\url{https://anonymous.4open.science/r/sci-tech-lag-C45D/}}

%% Bibliography
\begin{thebibliography}{9}

\bibitem{ahmadpoor2017dual}
M. Ahmadpoor, B. F. Jones,
The dual frontier: Patented inventions and prior scientific advance,
Science 357 (2017) 583--587.

\bibitem{marx2020reliance}
M. Marx, A. Fuegi,
Reliance on science: Worldwide front-page patent citations to scientific articles,
Strategic Management Journal 41 (2020) 1572--1594.

\bibitem{funk2017dynamic}
R. J. Funk, J. Owen-Smith,
A dynamic network measure of technological change,
Management Science 63 (2017) 791--817.

\bibitem{park2023papers}
M. Park, E. Leahey, R. J. Funk,
Papers and patents are becoming less disruptive over time,
Nature 613 (2023) 138--144.

\bibitem{priem2022openalex}
J. Priem, H. Piwowar, R. Orr,
OpenAlex: A fully-open index of scholarly works, authors, venues, institutions, and concepts,
arXiv preprint arXiv:2205.01833 (2022).

\bibitem{beltagy2019scibert}
I. Beltagy, K. Lo, A. Cohan,
SciBERT: A pretrained language model for scientific text,
in: Proceedings of EMNLP-IJCNLP, 2019, pp. 3615--3620.

\bibitem{ke2018comparing}
Q. Ke,
Comparing scientific and technological impact of biomedical research,
Journal of Informetrics 12 (2018) 706--717.

\bibitem{ke2020basicness}
Q. Ke,
Technological impact of biomedical research: The role of basicness and novelty,
Research Policy 49 (2020) 104071.

\bibitem{vaneck2010vosviewer}
N. J. van Eck, L. Waltman,
Software survey: VOSviewer, a computer program for bibliometric mapping,
Scientometrics 84 (2010) 523--538.

\bibitem{denter2025novelty}
N. M. Denter, J. Waterstraat, M. G. Moehrle,
Avoiding the pitfalls of direct linkage: A novelty-driven approach to measuring scientific impact on patents,
Journal of Informetrics 19 (2025).

\end{thebibliography}

\end{document}
